
\documentclass[11pt]{scrartcl}
\usepackage[sexy]{evan}

\title{Contest Collections}
\author{Guilherme}
\date{27 Octuber 2026}

\begin{document}

\maketitle

\section{Walkthroughs Graph theory}

\begin{problem}[HMMT February 2016]
  Let $n >1$ be an odd integer. On an $n\times n$ chessboard the center square and four corners are deleted. We wish to
  group the remaining $n^2-5$ squares into $\frac{1}{2} (n^2 -5)$ pairs, such that the two squares in each pair intersect
  at exactly one point ($i.e.$ they are diagonally adjacent, sharing a single corner)
For which odd integers $n > 1$ is this possible?
\end{problem}

\textbf{\color{green!15!darkgray}{Walkthrough.}}
Here I will follow the walkthrough by evan chen
\begin{soln}

\noindent

\begin{center}
\begin{asy}

size(2cm);


dot((0.5,0.5), black);
dot((2.5,0.5), black);
dot((0.5,2.5), black);
dot((1.5,1.5), black);
dot((2.5,2.5), black);

int n=3;
for (int i = 0; i<=n; ++i)
{draw((i,0)--(i,n), black);
draw((0,i)--(n,i), black);}


draw((0.5,1.5)--(1.5,0.5), red);
draw((1.5,2.5)--(2.5,1.5), red);
\end{asy}
\end{center}

This representation look a situation of the problem. As shown, the case first $n=3$ can do.

\begin{walk}
\item Can one do $n=3$?
  Yes, its possible.

\item Can one do $n=5$?

\begin{minipage}{0.3\textwidth}
\begin{asy}
size(4cm);


dot((0.5,4.5), black);
dot((4.5,4.5), black);
dot((2.5,2.5), black);
dot((0.5,0.5), black);
dot((4.5,0.5), black);

dot((1.5,4.5), red);
dot((4.5,3.5), red);
dot((3.5,0.5), red);
dot((0.5,1.5), red);

int n=5;
for (int i = 0; i<=n; ++i)
{draw((i,0)--(i,n), black);
draw((0,i)--(n,i), black);}

draw((1.5,4.5)--(0.5,3.5), red);
draw((3.5,4.5)--(4.5,3.5), red);
draw((4.5,1.5)--(3.5,0.5), red);
draw((1.5,0.5)--(0.5,1.5), red);

\end{asy}
\end{minipage}
\begin{minipage}{0.3\textwidth}
\begin{asy}
size(4cm);

dot((0.5,4.5), black);
dot((4.5,4.5), black);
dot((2.5,2.5), black);
dot((0.5,0.5), black);
dot((4.5,0.5), black);

fill(box((1,3), (2,4)), gray);
fill(box((3,3), (4,4)), gray);
fill(box((3,1), (4,2)), gray);
fill(box((1,1), (2,2)), gray);

dot((1.5,4.5), red);
dot((4.5,3.5), red);
dot((3.5,0.5), red);
dot((0.5,1.5), red);

dot((2.5,4.5), red);
dot((4.5,2.5), red);
dot((2.5,0.5), red);
dot((0.5,2.5), red);


int n=5;
for (int i = 0; i<=n; ++i)
{draw((i,0)--(i,n), black);
draw((0,i)--(n,i), black);}

draw((1.5,4.5)--(0.5,3.5), red);
draw((3.5,4.5)--(4.5,3.5), red);
draw((4.5,1.5)--(3.5,0.5), red);
draw((1.5,0.5)--(0.5,1.5), red);

draw((2.5,4.5)--(1.5,3.5), red);
draw((3.5,3.5)--(4.5,2.5), red);
draw((3.5,1.5)--(2.5,0.5), red);
draw((1.5,1.5)--(0.5,2.5), red);




\end{asy}
\end{minipage}

\par Please note that we are back to the case $n=3$.

\begin{center}
\begin{asy}
size(3cm);


dot((1.5,1.5), black);

dot((1.5,2.5), red);
dot((1.5,0.5), red);

fill(box((0,0), (1,1)), gray);
fill(box((2,0), (3,1)), gray);
fill(box((2,2), (3,3)), gray);
fill(box((0,2), (1,3)), gray);

int n=3;
for (int i = 0; i<=n; ++i)
{draw((i,0)--(i,n), black);
draw((0,i)--(n,i), black);}


draw((1.5,2.5)--(0.5,1.5), red);
draw((2.5,1.5)--(1.5,0.5), red);

\end{asy}
\end{center}

\paragraph{My point of view} This moment my strategies is connections of diagonally adjacent parallel to the center.
I think that is possible say all cases can do return back to the $n=3$ case.

\item Can one do $n=7$?
\par I try many times and can not do. I will get rid of it to draw a $7\times 7$.

\item Find a coloring of the squares by black and white so that diagonally adjacent squares are opposite colors.

The representations with balls ever parallel to the center.

\begin{center}
\begin{asy}
import patterns;
add("hach", hatch(1.5mm, NE, black+0.4));
size(3cm);

fill(box((1,4), (2,5)), pattern("hach"));
fill(box((2,4), (3,5)), pattern("hach"));
fill(box((4,3), (5,4)), pattern("hach"));
fill(box((4,2), (5,3)), pattern("hach"));
fill(box((3,0), (4,1)), pattern("hach"));
fill(box((2,0), (3,1)), pattern("hach"));
fill(box((0,1), (1,2)), pattern("hach"));
fill(box((0,2), (1,3)), pattern("hach"));

dot((0.5,4.5), black);
dot((4.5,4.5), black);
dot((2.5,2.5), black);
dot((0.5,0.5), black);
dot((4.5,0.5), black);

int n=5;
for (int i = 0; i<=n; ++i)
{draw((i,0)--(i,n), black);
draw((0,i)--(n,i), black);}
\end{asy}
\end{center}

\par The squares that to be left over in the center may have any fillers.

\item Use this to narrow down the set of possible n to two value
\par Why make to narrow down the set for two possible?


\begin{center}
\begin{asy}
import patterns;
add("hach", hatch(1.5mm, NE, black+0.4));
size(6cm);

fill(box((1,6), (2,7)), pattern("hach"));
fill(box((2,6), (3,7)), pattern("hach"));
fill(box((6,5), (7,6)), pattern("hach"));
fill(box((6,4), (7,5)), pattern("hach"));
fill(box((5,0), (6,1)), pattern("hach"));
fill(box((4,0), (5,1)), pattern("hach"));
fill(box((0,1), (1,2)), pattern("hach"));
fill(box((0,2), (1,3)), pattern("hach"));

fill(box((4,6), (5,7)), pattern("hach"));
fill(box((6,2), (7,3)), pattern("hach"));
fill(box((2,0), (3,1)), pattern("hach"));
fill(box((0,4), (1,5)), pattern("hach"));
dot((0.5,6.5), black);
dot((6.5,6.5), black);
dot((6.5,0.5), black);
dot((0.5,0.5), black);
dot((3.5,3.5), black);

int n=7;
for (int i = 0; i<=n; ++i)
{draw((i,0)--(i,n), black);
draw((0,i)--(n,i), black);}

draw((1.5,6.5)--(0.5,5.5), red);
draw((5.5,6.5)--(6.5,5.5), red);
draw((6.5,1.5)--(5.5,0.5), red);
draw((1.5,0.5)--(0.5,1.5), red);

draw((2.5,6.5)--(1.5,5.5), red);
draw((5.5,5.5)--(6.5,4.5), red);
draw((5.5,1.5)--(4.5,0.5), red);
draw((1.5,1.5)--(0.5,2.5), red);

draw((4.5,6.5)--(3.5,5.5), red);
draw((6.5,2.5)--(5.5,3.5), red);
draw((2.5,0.5)--(3.5,1.5), red);
draw((0.5,4.5)--(1.5,3.5), red);

label("$\cdots$", (3.5,0.5), black);
label("$\cdots$", (3.5,6.5), black);
label("$\vdots$", (6.5,3.5), black);
label("$\vdots$", (0.5,3.5), black);

\end{asy}
\end{center}


\paragraph{By Induction.} Repeating the algorithm, let's  to arrive in the case $n=3$ with different conditions.

\noindent

\begin{center}
\begin{asy}
size(2cm);

dot((0.5,1.5), orange);
dot((1.5,0.5), orange);
dot((1.5,1.5), black);
dot((1.5,2.5), orange);
dot((2.5,1.5), orange);

int n=3;
for (int i = 0; i<=n; ++i)
{draw((i,0)--(i,n), black);
draw((0,i)--(n,i), black);}

\end{asy}
\end{center}

\par With we're may, don't are possible for any $n$ odd integer $n>5$.

\item Wrap up the problem using your earlier work.

\par Therefore conclusions that $n=3$ and $n=5$ are units solutions.

\end{walk}

\end{soln}

\begin{problem}[\href{https://artofproblemsolving.com/community/c5h1823554p12189456}{USAJMO 2019}]
There are $a+b$ bowls arranged in a row, numbered $1$ through $a+b$, where $a$ and $b$ are given positive integers.
Initially, each of the first $a$ bowls contains an apple, and each of the last $b$ bowls contains a pear. A legal move
consists of moving an apple from bowl $i$ to bowl $i+1$ and a pear from bowl $j$ to bowl $j-1$, provided that the
different $i-j$ is even. We permit multiple fruits in the same bowl at the same time. The goal is to end up with the
first $b$ bowls each containing a pear and the last $a$ bowls each containing an apple. Show that this is possible if
and only if the product $ab$ is even.
\end{problem}

\paragraph{My solution.} This is problem before see the walkthrough of the Evan. I got it solving the problem in little
plus of 1 hour, but don't know how to put it here. Hence, I put it here solution by evan with my notes.\par
Look the walkthrough, I see that my solution is wasn't entirely accurate.

\begin{soln}
\item\textbf{\color{green!15!darkgray}{Walkthrough.}}  First we show that if $ab$ is even then the goal is possible. We prove the result by induction on $a + b$.

\begin{remark}
  This idea is very very useful. I really don't would think this strategies but, it is good. Verify a thing that is
  possible could help to look because it works.
\end{remark}

\begin{walk}
\item If $\min(a,b) \ge 1$, and $a$ and $b$ are opposite parity, show that in one swap one can reduce from $(a,b)$
to $(a-1,b-1)$.

\item If $\min (a,b) \ge 2$, and $a$ and $b$ are both even, show that in two swaps one can reduce from $(a,b)$ to
$(a-2,b-2)$.

\item Formulate a set of base cases and complete the proof via induction.
\end{walk}
Conversely, we need to show that the task is impossible when $ab$ is odd.

\begin{walk}[start=4]
\item Let $X$ denote the number of apples in odd-numbered bowls, and $Y$  denote the number of pears ind odd-numbered
bowls. Find a relation between $X$ and $Y$ that does not change under the operation.

\item Use this to show that the task is impossible when $ab$ is odd.
\end{walk}

\item\textbf{\color{red!40!darkgray}{Solution.}}

\begin{itemize}
\item If $\min(a,b) =0$ there is nothing to check.

\item If $\min (a,b)=1$, say $a=1$, then $b$ is even  and we can swap the (only) leftmost apple with the rightmost
pear by working only with those fruits.

\begin{remark}
In fact, we can repeat the process trading $m - p - p$, whenever have one even-numbers of pears to right of one apple.
\end{remark}


\item Now assume $\min(a, b) \ge 2$ and $a + b$ is odd. Then we can swap the leftmost apple with rightmost pear by working
only with those fruits, reducing to the situation of $(a-1, b-1)$ which is possible by induction (at least one of
them is even).

\begin{remark}
Say $a$ is odd and $b$ is even.\par
We have the configuration: $\{m,m,\ldots,m,p,\ldots,p,p\}$. We may trade the leftmost apple with rightmost pear, $i=1$
and $j=$odd, $j-i$ will be even. With this, reducing to the situation for $(a-1,b-1)$, which is possible by induction.
\end{remark}

\item Finally, assume $\min(a,b) \ge 2$ and $a+b$ is even ($i.e$ $a$ and $b$ are both even). Then we can swap the apple in
position $1$ with the pear in position $a+b-1$, and apple in position $2$ with the pear in the position $a+b$.
This reduce to the situation of $(a-2,b-2)$ which is also possible by induction.

\begin{remark}
I don't many thing to talk for the item, I think that this feature is one of the but I want for me. See that in the
problem \href{https://artofproblemsolving.com/community/c6h3181582p28989993}{OBM 2023/1}. The solution of the Luciano
Castro by \href{https://www.youtube.com/watch?v=rPlnud-oexQ&t=4s}{Matematicamente} is very resilient. I will study
to the solution.


\end{remark}
\end{itemize}

Now we show that the result is impossible if $ab$ is odd. Define\par
\begin{center}
$X=$number apples in odd-numbered bowls\\
$Y=$number pears in odd-numbered bowls.
\end{center}

Note that $X-Y$ does not change under this operation. However, if $a$ and $b$  are odd, then we initially have
$X=\frac{1}{2}(a+1)$ and $Y=\frac{1}{2}(b-1)$, while the target position has $X=\frac{1}{2}(a-1)$ and
$Y=\frac{1}{2}(b+1)$. So, when $ab$ id odd this is not possible.

\end{soln}

\begin{remark}[Evan]
Remark. Another proof that $ab$ must be even is as follows.\par
First, note that apples only move right and pears only move left, a successful operation must take exactly $ab$ moves.
So it is enough to prove that the $number$ $of$ $moves$ made must be even.\par
However, the number of fruits in odd-numbered
bowls either increases by $+2$ or $-2$ in each move (according to whether $i$ and $j$ are both even or both odd), and
since it ends up being the same at the end, the number of move must be even.\par
Alternatively, as pointed out in the official solutions, one can consider the
sums of squares of positions of fruits. The quantity changes by\par

\begin{align*}
  [(i+1)^2+(j-1)^2] - (i^2+j^2) = 2(i-j) + 2 \equiv 2 \pmod{4}
\end{align*}

at each step, and eventually the sums of squares returns to zero, as needed.
\end{remark}

\paragraph{Double Counting} The form that he used by conclusion the solution was very ingenious, a double counting is
one walk to think in this idea. The normal counting is $X=\frac{1}{2}(a+1)$ since $a$ is odd, hence, there will be more
odd-numbered bowls than even-numbered bowls. $Y=\frac{1}{2}(b-1)$ since $a+1$ is even, therefore, a counting of bowls
with pears starting in even bowl and ends also in even bowls. \par
``{\textbf{While the target position}} ", why will be this target position?\par
The target position in my point of view is when swap the leftmost apple with the rightmost pear. The apple in rightmost
position $a+b$ is even and pear in leftmost position $1$ is odd, as soon as the swap is done, $X-1=\frac{1}{2}(a-1)$
and $Y+1=\frac{1}{2}(a+1)$. What changes everything.\par
I'll speak again, beautiful idea to ends of the problem.


\end{document}
